\section{Síntesis de los principales hallazgos}

La investigación permitió implementar un prototipo que integra procesamiento de datos crediticios, entrenamiento automatizado, registro de experimentos, almacenamiento analítico y visualización. En la evaluación presentada, LightGBM mostró una mayor capacidad discriminativa que las implementaciones de CNN y MLP. El rendimiento de LightGBM disminuyó conforme aumentó el horizonte, y el recall presentó una reducción especialmente relevante en los plazos largos.

El análisis del código fuente permitió documentar completamente las ocho condiciones de \texttt{crisis\_flag, las 21 características, la arquitectura exacta de cada modelo y el protocolo de división de datos (49\% entrenamiento, 21\% validación, 30\% prueba).}

\section{Respuesta al objetivo general}

El objetivo general se alcanzó en lo relativo al diseño y la integración técnica de los componentes principales. Se implementaron flujos de entrenamiento e inferencia en Apache Airflow, registro en MLflow, almacenamiento en PostgreSQL y dashboards en Apache Superset. El resultado debe presentarse como un prototipo funcional. La evidencia disponible no permite concluir que la plataforma controle la mora, reduzca pérdidas o haya sido validada para apoyar decisiones en un entorno productivo.

\section{Respuesta a los objetivos específicos}

Respecto al primer objetivo, se identificaron factores plausibles asociados con el comportamiento de las redes neuronales: una muestra secuencial reducida (solo bloques con $\geq 24$ meses de historial generan secuencias), desbalance de clases sin compensación en MLP, ventanas de seis meses y configuración documentada completa. El MLP no implementa \texttt{sample\_weights, lo que puede explicar su tendencia a predecir todas las observaciones como ``no crisis''. No se realizaron experimentos de ablación, por lo que estos factores no pueden declararse causas demostradas.}

En relación con el segundo objetivo, LightGBM alcanzó el mayor AUC-ROC bajo las configuraciones disponibles. Esta comparación no demuestra que los modelos basados en árboles sean superiores en todos los conjuntos tabulares, porque la búsqueda de hiperparámetros no fue equivalente (LightGBM usa semilla fija y tuning documentado; CNN y MLP no fijan semilla).

Para el tercer objetivo, se observó una disminución del rendimiento de LightGBM con el horizonte. El AUC-ROC pasó de 0,874 a un mes a 0,702 a dieciocho meses, mientras que el recall descendió de 0,542 a 0,034. Por tanto, los horizontes largos deben interpretarse con cautela y no pueden considerarse operativamente útiles sin calibración, selección de umbral y análisis de costos de error.

El cuarto objetivo se cumplió como implementación técnica. El datamart integra dimensiones de tiempo, riesgo, sector y sucursal, y los dashboards permiten consultar información histórica y predicciones. No se evaluó la utilidad con usuarios, la latencia ni el impacto sobre decisiones, por lo que esas dimensiones permanecen pendientes.

\section{Contribuciones}

La primera contribución es una arquitectura de extremo a extremo que enlaza ETL, modelado, MLOps, datamart y visualización. La segunda es la comparación de tres enfoques para dieciocho horizontes, junto con el análisis de la degradación temporal. La tercera es la implementación de un esquema estrella y dashboards que integran predicciones y agregados históricos.

Las contribuciones se circunscriben al prototipo y a la institución analizada. No se afirma superioridad universal de LightGBM ni se generalizan los resultados a otras entidades financieras.

\section{Limitaciones}

El estudio presenta limitaciones metodológicas y de generalización. La definición de \texttt{crisis\_flag quedó documentada mediante ocho condiciones con pesos acumulativos (umbral $\geq 5$), pero existe un riesgo confirmado de target leakage: las variables judiciales (\texttt{tasa\_judicial}, \texttt{creditos\_judiciales}, \texttt{total\_costo\_judicial}) se utilizan tanto para definir la etiqueta como como características de entrada. Esta dependencia puede inflar artificialmente el poder predictivo de esas variables y del modelo.}

La división de datos sigue un protocolo temporal con tres conjuntos (49\% entrenamiento, 21\% validación, 30\% prueba), y el early stopping utiliza exclusivamente el conjunto de validación. Sin embargo, no se reportan conteos por horizonte, no se fijan semillas para CNN y MLP, y no se realizan repeticiones para estimar variabilidad.

No se reportan PR-AUC, calibración, curvas por umbral ni análisis detallado de falsos negativos. La comparación entre modelos utiliza presupuestos de ajuste no equivalentes. Además, el estudio se realizó en una sola institución y no incluye validación externa.

Desde la perspectiva operativa, las capturas demuestran la implementación de los flujos y dashboards, pero no se presentan pruebas de rendimiento, disponibilidad, recuperación ante fallos, integridad de datos o usabilidad. Estas limitaciones impiden considerar el prototipo como un sistema productivo validado.

\section{Trabajo futuro}

El trabajo futuro debe priorizar la reproducibilidad. Se requiere registrar conteos por etapa y horizonte, fijar semillas para todos los modelos, realizar repeticiones con múltiples semillas y versionar datos, código y entornos.

Es crítica una ablación de leakage: repetir la evaluación excluyendo las variables judiciales (\texttt{tasa\_judicial, \texttt{creditos\_judiciales}, \texttt{total\_costo\_judicial}) para estimar cuánto del poder predictivo se debe a la dependencia con la etiqueta.}

La comparación de modelos debe equilibrar los presupuestos de ajuste, utilizar múltiples semillas, reportar intervalos de confianza, incluir PR-AUC, calibración y selección de umbrales basada en costos. Para las redes neuronales, se recomienda añadir \texttt{sample\_weights al MLP y explorar arquitecturas con mayor capacidad.}

Finalmente, el prototipo debe evaluarse con datos de otras instituciones y mediante pruebas funcionales y de usabilidad. Solo después de estas validaciones sería posible estimar su contribución a procesos de decisión, alertas o gestión de cartera.
